\documentclass[11pt]{article}
\usepackage[margin=1in]{geometry}
\usepackage{amsmath,amssymb,amsthm}
\usepackage[hidelinks]{hyperref}
\newtheorem{theorem}{Theorem}
\title{A nine-step recurrence for the decimal chunks of \(12321\)}
\author{Pioneer project}
\date{28 November 2026}
\begin{document}
\maketitle

\begin{abstract}
Moving nine chunks forward in the periodic word \(12321\,12321\cdots\) appends one fixed block of five digits to each chunk. This proves the recurrence and rational generating function conjectured for OEIS A321225.
\end{abstract}

Partition the infinite word \(12321\,12321\cdots\) into consecutive chunks whose digit sums are the successive positive integers, and read the \(n\)-th chunk as a decimal integer \(a_n\). The recurrence and generating function below were conjectured by Chai Wah Wu in the OEIS entry \cite{oeis}.

\begin{theorem}
The chunks exist uniquely, and for \(n>18\),
\[
a_n=100001a_{n-9}-100000a_{n-18}.
\]
The ordinary generating function is
\[
\sum_{n\ge1}a_nx^n
=\frac{P(x)}{(1-x^9)(1-100000x^9)},
\]
where
\begin{align*}
P(x)={}&x+2x^2+3x^3+211x^4+23x^5+2112x^6+3211x^7+2321x^8\\
&+12321x^9+23210x^{10}+32110x^{11}+21120x^{12}+23000x^{13}\\
&+21100x^{14}+30000x^{15}+20000x^{16}+10000x^{17}.
\end{align*}
\end{theorem}

\begin{proof}
Let \(S(j)\) be the digit sum of the first \(j\) digits and \(T_n=n(n+1)/2\). A cut after the \(n\)-th chunk must satisfy \(S(c_n)=T_n\). If \(j=5q+r\) with \(0\le r<5\), then
\[
S(j)=9q+u_r,\qquad (u_0,u_1,u_2,u_3,u_4)=(0,1,3,6,8).
\]
Since \(T_{n+9}=T_n+9(n+5)\), the residues of \(T_n\) modulo 9 repeat with period 9. For \(n=0,1,\ldots,8\) they are \(0,1,3,6,1,6,3,1,0\). Each occurs among the \(u_r\), so each \(c_n\) exists. The digits are positive, so \(S\) is strictly increasing and each cut is unique. In particular \(c_0=0\), and \(c_n>c_{n-1}\) for \(n\ge1\).

The cut formula and triangle identity show that
\[
c_{n+9}=c_n+5(n+5).
\]
This also gives the periodic increase in the chunk lengths listed in A321232 \cite{lengths}. Both ends of the chunk at \(n+9\) have the same residues modulo 5 as those of the chunk at \(n\). The right end moves five more places than the left end, so the new chunk agrees with the old one up to the old length and then has five extra digits. Let \(B_n\) be the decimal integer made from the five digits immediately after \(c_n\). Then
\[
a_{n+9}=10^5a_n+B_n,\qquad B_{n+9}=B_n.
\]
Eliminating \(B_n\) between the identities at \(n\) and \(n+9\) gives the recurrence.

For the numerator, let \(a_0=0\) and give negative indices coefficient zero. After multiplication by \((1-x^9)(1-100000x^9)\), the coefficients through degree 9 are \(a_0,\ldots,a_9\). At degree \(n+9\) with \(1\le n\le9\), the coefficient is \(B_n-a_n\). All subsequent coefficients vanish by the recurrence. The cut formula gives the complete list of required initial values and blocks:
\[
\begin{array}{c|r|r@{\qquad}c|r|r@{\qquad}c|r|r}
n&a_n&B_n&n&a_n&B_n&n&a_n&B_n\\ \hline
1&1&23211&4&211&23211&7&3211&23211\\
2&2&32112&5&23&21123&8&2321&12321\\
3&3&21123&6&2112&32112&9&12321&12321
\end{array}
\]
Substitution gives \(P(x)\). For example, the block after the first chunk is \(23211\), so \(a_{10}=10^5+23211=123211\).
\end{proof}

\begin{thebibliography}{2}
\bibitem{oeis}
OEIS Foundation Inc., entry A321225, submitted by Seiichi Manyama (31 October 2018), with conjectures by Chai Wah Wu (18 April 2024).
\url{https://oeis.org/A321225}
\bibitem{lengths}
OEIS Foundation Inc., entry A321232, lengths of the decimal chunks.
\url{https://oeis.org/A321232}
\end{thebibliography}
\end{document}
